\section{Sistem Informasi Manajemen (SIM)}
\label{sec:sim}

Sistem Informasi Manajemen adalah serangkaian komponen yang saling terkait yang bertujuan untuk mengumpulkan, memproses, menyimpan, dan mendistribusikan informasi untuk mendukung pengambilan keputusan, koordinasi, dan pengendalian dalam suatu organisasi \cite{armah2024sim}. Sistem Informasi Manajemen memiliki tujuan utama untuk mendukung kegiatan operasional dan strategi organisasi dengan cara memberikan informasi yang tepat dan akurat kepada manajemen. Dengan demikian, SIM dapat membantu manajemen dalam mengambil keputusan yang lebih baik, meningkatkan efisiensi dan efektivitas operasional, serta meningkatkan kualitas pengelolaan administrasi. 

Sistem Informasi Manajemen terdiri atas beberapa komponen yang meliputi komponen input, komponen model, komponen output, komponen teknologi, komponen hardware, komponen software, komponen basis data, dan komponen kontrol \cite{rusdiana2014sim}. Seluruh komponen tersebut saling berinteraksi satu sama lain untuk mencapai tujuan sistem informasi manajemen. Penjelasan setiap komponen nya adalah sebagai berikut:

\begin{itemize}
    \item Komponen Input \\
    Input mewakili data yang masuk ke dalam sistem informasi. Input termasuk metode dan media untuk menangkap data yang akan dimasukkan, yang dapat berupa dokumen-dokumen dasar.
    \item Komponen Model \\
    Komponen ini terdiri atas kombinasi prosedur, logika, dan model matematik yang akan memanipulasi data input dan data yang tersimpan di basis data dengan cara yang telah ditentukan untuk menghasilkan keluaran yang diinginkan.
    \item Komponen Output \\
    Hasil dari sistem informasi adalah keluaran informasi yang berkualitas dan dokumentasi yang berguna untuk semua pemakai sistem.
    \item Komponen Teknologi \\
    Teknologi merupakan "tool box" dalam sistem informasi. Teknologi digunakan untuk menerima input, menjalankan model, menyimpan dan mengakses data, menghasilkan dan mengirimkan keluaran, dan membantu pengendalian dari sistem secara keseluruhan.
    \item Komponen Hardware \\
    Hardware berperan penting sebagai suatu media penyimpanan vital bagi sistem informasi yang berfungsi sebagai database sumber data dan informasi untuk memperlancar dan mempermudah kerja dari sistem informasi.
    \item Komponen Software \\
    Software berfungsi sebagai tempat untuk mengolah, menghitung, dan memanipulasi data yang diambil dari hardware untuk menciptakan suatu informasi.
    \item Komponen Basis Data \\
    Basis data (database) merupakan kumpulan data yang saling berkaitan dan berhubungan satu dengan yang lain, tersimpan di perangkat keras komputer dan menggunakan perangkat lunak untuk Sistem Informasi Manajemen memanipulasinya. Data perlu disimpan dalam basis data untuk keperluan penyediaan informasi lebih lanjut. Data dalam basis data perlu diorganisasikan sedemikian rupa informasi yang dihasilkan berkualitas. Organisasi basis data yang baik juga berguna untuk efisiensi kapasitas penyimpanannya. Basis data diakses atau dimanipulasi menggunakan perangkat lunak yang disebut Database Management System (DBMS).
    \item Komponen Kontrol \\
    Banyak hal yang dapat merusak sistem informasi, seperti bencana alam, kegagalan sistem, ketidakefisienan, sabotase, dan sebagainya. Beberapa pengendalian perlu dirancang dan diterapkan untuk meyakinkan bahwa hal yang dapat merusak sistem dapat dicegah ataupun jika telanjur terjadi kesalahan-kesalahan dapat langsung cepat diatasi.
\end{itemize}

Dalam konteks proyek ini, aplikasi Monitoring Pajak Daerah (MonPD Reborn) dikembangkan sebagai wujud dari Sistem Informasi Manajemen. Aplikasi ini berfungsi untuk mengolah data-data transaksi dan realisasi pajak daerah yang bersifat mentah menjadi informasi strategis. Informasi tersebut disajikan dalam bentuk dashboard, rekapitulasi, dan laporan yang membantu pimpinan dan staf di Bapenda Kota Surabaya dalam memantau kinerja penerimaan pajak, mengidentifikasi potensi tunggakan, dan membuat keputusan yang lebih baik terkait strategi optimalisasi pendapatan daerah.